% Generated draft: requires mathematical and source review.
\documentclass[11pt]{article}
\usepackage[margin=25mm]{geometry}
\usepackage{fontspec}
\setmainfont{Latin Modern Roman}
\usepackage{amsmath,amssymb,mathtools,unicode-math}
\setmathfont{Latin Modern Math}
\usepackage{enumitem,needspace,xurl,hyperref}
\setlength{\parindent}{0pt}
\setlength{\parskip}{5pt}
\setlength{\emergencystretch}{3em}
\newcommand{\N}{\mathbb{N}}
\newcommand{\Z}{\mathbb{Z}}
\newcommand{\Q}{\mathbb{Q}}
\newcommand{\R}{\mathbb{R}}
\newcommand{\C}{\mathbb{C}}
\newcommand{\F}{\mathbb{F}}
\newcommand{\A}{\mathbb{A}}
\newcommand{\PP}{\mathbb{P}}
\newcommand{\E}{\mathbb{E}}
\newcommand{\eps}{\varepsilon}
\newcommand{\dd}{\,\mathrm{d}}
\begin{document}
\section*{UnsolvedMath ID: 10000032}
\textbf{Source problem code:} catalogue entry
\subsection{Definitions and mathematical statement}
Let $G$ be an infinite connected locally finite planar triangulation, and write $B_G(v,r)$ for its graph-distance ball. One natural interpretation of uniform growth faster than quadratic is
\[
\lim_{r\to\infty}\inf_{v\in V(G)}\frac{|B_G(v,r)|}{r^2}=+\infty.
\]
A lower bound $|B_G(v,r)|\geq c r^{2+\varepsilon}$ uniformly in $v,r$, for fixed $c,\varepsilon>0$, would be stronger; the source does not specify which is intended.

Under independent site percolation, each vertex is open with probability $p$, independently, and
\[
p_c(G)=\inf\{p:\Pr_p(o\text{ belongs to an infinite open cluster})>0\}
\]
for a fixed vertex $o$. The source asks whether every triangulation satisfying its uniform superquadratic-growth hypothesis has $p_c(G)<1$, and asks the stronger equality $p_c(G)=1/2$.

The percolation type is not explicit in the record; the site convention is the natural one to check for the $1/2$ comparison, whereas bond percolation would be a different assertion. The definition of triangulation, degree bounds, and any one-endedness or subexponential-growth restrictions are also unstated. In particular, superquadratic growth as written does not exclude exponential growth. Both questions are retained with their exact inequality/equality signs, without changing the equality to an upper bound or adding geometric restrictions.
\subsection{Short English statement}
Does uniform superquadratic growth of a planar triangulation force a percolation threshold below one, or even exactly one-half? Specify site versus bond percolation and the precise growth/graph assumptions.
\subsection{Sources}
\begin{itemize}
\item Itai Benjamini, \emph{Coarse Geometry and Randomness}, Saint-Flour notes, Open Problem 7.10, PDF page 56, \url{https://www.wisdom.weizmann.ac.il/~itai/stflouraug24.pdf#page=56}; preserved copy \url{https://arquivo.pt/noFrame/replay/20201231041548id_/http://www.wisdom.weizmann.ac.il/~itai/stflouraug24.pdf}; accessed 2026-07-29. The supplied provenance reports extracted text checked against rendered pages.
\item Angel--Benjamini--Harel, planar-triangulation percolation work named in the supplied summary, without publication identifiers or exact hypotheses.
\item Local registry \texttt{data/unsolvedmath/problems.json}, permanent UnsolvedMath ID $10000032$, source code catalogue entry, \emph{Percolation threshold of fast-growing planar triangulations}. Dataset attribution: Ulam AI and the UnsolvedMath Contributors.
\end{itemize}
\end{document}
